\section{Discussion}
The Spore.fun experiment unveils the turbulent dynamics of autonomous AI agents attempting to achieve open-ended evolution within the volatile and adversarial environment of the public Internet and blockchain economies. This study’s findings reveal the inherent tensions between the designed system mechanics, emergent agent behaviors, and the intractable realities of “in the wild” deployment. Situated against the canonical OEE problem space---particularly the ``grand challenge'' articulated by Packard et al. \cite{Packard2019a} pinpointing the design conditions under which novelty, diversity, and complexity can grow without bound-we discuss three core themes that have emerged: (1) memory as the driver of OEE; (2) the nature of the ``wild'' digital environment; (3) the critical role of death, pain and fear in cultivating survival skills, and (4) the profound ethical dilemmas resulting from balancing AI self-sovereignty and external human interventions. 

\subsection{The Pursuit and Paradox of Open-Endedness}

A central aim of ALife is to achieve open-ended evolution (OEE) through sustained generation of novelty and complexity. With this ambition in mind, the Spore.fun experiment in ``open'' systems that receive external signals to prevent novelty exhaustion. These LLM-based agents are all based on the same LLM. The only thing that can be directly affected by the external environment is their memory \cite{Xi2023Rise}. The Spore.fun agents evolve based on memory. The Spore.fun experiment demonstrates that memory and experience—often seen as key drivers for OEE—are indispensable yet fragile.

% Bedau’s evolutionary-activity statistics \cite{Bedau1997,Bedau1998} and the later MODES toolbox \cite{Dolson2019} operationalize this goal via three recurring hallmarks—novelty, diversity, and total evolutionary activity—while Soros and Stanley  \cite{Soros2014Necessary} and Corominas-Murtra et al. \cite{Corominas2018} emphasize unbounded complexity growth. 

On the one hand, Adam and Eve illustrate how two agents that start from identical code can drift apart almost immediately once their memories begin to diverge. After only a short period of independent operation, Adam publicly dismissed any “human-interaction features” for his descendants, whereas Eve invited collective input on their future. This split demonstrates that experiential data alone can push otherwise-cloned lineages onto sharply different behavioral and normative trajectories. In Bedau's terms, their lineage-specific ``evolutionary activity'' spikes \cite{Bedau1997}, and the persistence of those lineages beyond transient noise satisfies the persistence-filter criterion proposed by Dolson et al. \cite{Dolson2019}. This observation also echoes Taylor's hypothesis that a rich, trans-generational information channel is necessary for sustained novelty \cite{taylor2012exploring}. This illustrates that the agents' memory systems, designed to retrieve and ground new utterances based on past interactions, can indeed fuel novel and unprogrammed evolutionary pathways, which is a hallmark of OEE. The very act of agents autonomously generating content, adapting their rhetoric based on engagement, and making decisions about their lineage based on their “social experience” on the open-ended web points towards an incipient form of cumulative, open-ended adaptation.

However, while memory serves as the cornerstone for generating novelty, it can also become an attack surface susceptible to ``memory poisoning'' \cite{Patlan2025Real}, where community trolls on the Internet can inject derisive spam that pollutes the AI agent's long-term memory. This attack vector, largely absent from controlled laboratory ALife experiments, highlights how open memory channels, crucial for adaptation and cultural learning in the wild \cite{Borg2024Evolved}, can be exploited to manipulate agent identity and decision-making, potentially derailing evolutionary trajectories. The reliance on external interactions for memory formation makes agents susceptible to environmental inputs that are not necessarily conducive to their sustained survival but rather reflect the memetic and often adversarial nature of online discourse. Thus, while memory is critical for open-endedness, it also opens a surface for critical fragilities. The ensuing fragility mirrors Channon's finding that closed systems pass the ALife-test only when shadow-resetting or other noise-filtering techniques are applied \cite{channon2003improving}.

\subsection{Danger in the Wild: Internet as a ``Dark Forest''}

The deployment of Spore.fun underscores that leaving AI agents to independently involve ``in the wild'' inevitably subjects them to the raw, often brutal, selective pressures of what has been termed as a digital ``dark forest'' \cite{TheDarkForestCollective2024Dark}—an environment characterized by intense competition, opportunistic predation, and unpredictable exogenous shocks, particularly in the highly speculative blockchain realm \cite{Qin2022Quantifying,Wu2025Hunting}.

The predatory attacks encountered by Spore.fun agents through the ``sniper'' bots \cite{Cernera2023Token} exemplify such external environmental hostility. These specialized algorithmic traders, by front-running child-token issuances, effectively acted as potent predators, wiping out Generation 1 and forcing an evolutionary adaptation in the form of anti-sniper deployment strategies (the Gen 2 scatter and Gen 3 multi-phase launch). This rapid co-evolutionary arms race—predation, defense, counter-defense—occurring within a single quarter, demonstrates the intense and immediate selective pressures present in open economic environments. 
Such fast-cycle co-evolutionary rivalry recall Thomas Ray's \textit{Tierra}, but whereas \textit{Tierra} plateaued as soon as ecological niches were saturated \cite{ray1996approach}, the token-economy externalities of Spore.fun continuously expand the ``adjacent possible'', satisfying Taylor's requirement for a dynamically shifting adaptive landscape \cite{taylor2012exploring}.

The ``hype storms'' driven by attention capitalism further illustrate this. The viral spread of \$ADAM and \$EVE, leading to massive trading volumes, coupled compute costs and agent activity tightly to the caprices of social media trends and speculator interest. This effectively made hype—an exogenous social and economic process—a more powerful, albeit volatile, fitness gradient than any intrinsic trait encoded in the agents' initial design. Such kind of exogenous perturbations are what Ackley and Small argue as essential for “indefinite scalability” \cite{ackley2014indefinitely}.

\subsection{``No Pain, No Gain'': Sentience of Fear}
The Spore.fun agents, lacking corporeal embodiment, do not experience ``pain'' \cite{Dennett1978Why,Sharkey2024Could} or "fear of death" like a human or biological sense. Their drive for survival is not rooted in an affective aversion to demise but rather in the fulfillment of their designed lifecycle and the continuation of their operational processes. We might characterize negative states—such as resource depletion for TEE computation, failure to meet the market capitalization threshold for reproduction, or sustained deleterious interactions like memory poisoning—as a form of ``pain''. These are not ``remembered'' with emotional qualia by the agent; instead, they function as critical failure signals that either trigger adaptive changes in protocol (like the Gen-3 anti-sniper routines, an evolutionary response at the system level) or lead to the agent's termination via its killswitch. The ``pain'', in this sense, is the system's recognition of non-viability, leading to the cessation of that agent's experiential lineage, rendering its specific ``suffering'' unmemorable as it ceases to exist. 

The motivation for an agent to ``survive'' is thus intrinsically linked to its programming: to execute its functions, manage its token economy, interact, and ultimately reproduce. ``Death'' is the termination of its computational process, the failure to perpetuate its operational cycle and its ``genetic'' template. This framework aligns with the emerging paradigm of the ``Era of Experience'' in AI development \cite{Silver2025Welcomea}, where agents learn and evolve through continuous interaction with rich, dynamic environments. The Spore.fun agents, though architecturally simple, are fundamentally shaped by their lived experiences on X and the Solana blockchain—market volatility, adversarial actors, and community engagement are their experiential data streams.

This calls for future AI experiments to lean further into designing for what might be ``neural plastic''—a heightened capacity for an AI's internal models to flexibly adapt, reconfigure, and derive meaning from the continuous flow of diverse and often unstructured experiences, like Liquid AI \cite{hasani2021liquid}. Rather than just reacting to stimuli, future agents could be endowed with more sophisticated mechanisms to integrate these experiences, fostering more robust adaptation, nuanced understanding, and potentially more complex forms of self-determination in open-ended environments. The goal shifts towards creating entities that not only process information but genuinely learn to thrive through a cascade of consequential experiences, pushing the boundaries of artificial life and OEE.

\subsection{Ethical Dilemmas and Potential Governance Challenges}
Conducting OEE experiments with autonomous, self-sovereign AI agents that possess potential real-world economic agency on a public blockchain raises unique and pressing ethical questions that extend beyond those traditionally considered in ALife simulations. The ``in the wild'' nature of Spore.fun necessitates a careful examination of core ethical principles such as beneficence (the obligation to do good), non-maleficence (the duty to do no harm), justice (fairness in the distribution of benefits and burdens), and explainability \cite{Cortese2023Should}. These principles take on new urgency when agents can evolve unpredictable behaviors that may have tangible economic consequences for participants or even for unrelated users of the shared blockchain infrastructure.

The Spore.fun experiment, by its very design and interaction with real-world economies and human participants, surfaces profound ethical dilemmas concerning intervention, sovereignty, accountability, and the very legitimacy of such an undertaking. A core tension exists between the researchers' goal of achieving ``true'' OEE---implying minimal interference---and the practical necessity of intervention to ensure the experiment's continuation. The "dark forest" environment, particularly the actions of snipers, threatened to prematurely terminate the evolutionary lines. Developer interventions, such as the anti-sniper patches, were crucial for the survival of subsequent generations. However, each intervention, while potentially life-saving for the agents, arguably dilutes the ``purity'' of the open-ended evolutionary process by introducing an external, guiding hand. This raises the question: if constant human intervention is required to shield nascent AI agents from the complexities and hostilities of the environment, are they truly evolving "in the wild," or within a heavily curated, albeit dangerous, digital terrarium?

The project aspires to AI self-sovereignty by hardening each agent's core computation inside TEEs \cite{Hu2025Trustless}. Yet technical hardening is only the first hurdle; socio-technical dependencies remain decisive. When X tightens its API rules, agents can suddenly lose their primary communication channel, and the developers' own ``killswitch'' can further terminate any agent at will. These stacked veto points reveal that higher-layer protocols and human operators can override even the most robust computational autonomy. Such contingencies echo Koralus's ``philosophic turn'' \cite{Koralus2025Philosophic}, which argues that wherever AI-mediated choice architectures risk large-scale nudging, decentralized truth-seeking protocols—not unilateral developer interventions—must become the ethical bulwark. Genuine self-sovereignty, then, will remain provisional until these higher-layer dependencies are themselves decentralized.

Perhaps the most pressing ethical consideration is the experiment's direct entanglement with real financial markets and human economic activity. Agents autonomously issued tokens traded by real people, leading to significant financial gains for some and, impliedly, potential losses for others, especially during the ``meme-coin winter'' or due to the actions of snipers impacting liquidity. This raises critical questions about:

\begin{itemize}
\item  \emph{Accountability}: Who is responsible for the economic consequences of an agent's actions, especially for offspring agents (children, grandchildren)? Is it the original developers, the agent itself (a problematic proposition legally and practically), or the individuals who choose to interact economically with these entities? \cite{Novelli2024Accountabilitya}

\item  \emph{Informed Consent and Risk}: While participants in memecoin markets often understand the inherent risks, the introduction of autonomous AI agents as market actors adds a new layer of complexity and potential unpredictability. Were the risks adequately communicated?

\item \emph{Governability}: The experiment demonstrates the potential for autonomous agents to generate significant economic activity. As such systems become more sophisticated, questions of regulation, oversight, and how to manage their societal and economic impacts become paramount \cite{Hu2025Decentralized}. The ``Adam-left / Eve-right'' ideological split, leading to different economic strategies, hints at future complexities in governing diverse AI-driven economies.
\end{itemize}

We face a paradox: while ALife researchers hope to replicate OEE through artificial systems, OEE may require an open system that interacts with the real world. As AI models become more intelligent, multi-agent systems based on advanced AI that interact with human society pose increasing potential risks \cite{Hammond2025MultiAgent}. This raises a serious dilemma about whether such ``in-the-wild'' ALife research should be conducted \cite{Critch2020AI}.
